\documentclass[11pt,a4paper]{article}
\usepackage[margin=1.8cm]{geometry}
\usepackage{xcolor}
\usepackage{amssymb}
\usepackage[most]{tcolorbox}
\usepackage{array,tabularx,booktabs}
\usepackage{fontspec}
\usepackage{polyglossia}
\setdefaultlanguage{french}
\setotherlanguage{arabic}
\setmainfont{DejaVu Serif}
\newfontfamily\arabicfont[Script=Arabic,Scale=1.05]{DejaVu Sans}
\newfontfamily\arabicfontsf[Script=Arabic,Scale=1.05]{DejaVu Sans}
\setlength{\parindent}{0pt}
\setlength{\parskip}{4pt}
\renewcommand{\arraystretch}{1.35}

\definecolor{bleu}{HTML}{2E5C9A}
\definecolor{vert}{HTML}{2E7D32}
\definecolor{rouge}{HTML}{C62828}
\definecolor{orange}{HTML}{E65100}

\newtcolorbox{carte}{colback=bleu!6,colframe=bleu,title=Où on est sur la carte,fonttitle=\bfseries}
\newtcolorbox{examen}{colback=vert!6,colframe=vert,title=Phrases d'examen à connaître (en français),fonttitle=\bfseries}
\newtcolorbox{piege}{colback=rouge!5,colframe=rouge,title=Pièges des examens,fonttitle=\bfseries}
\newtcolorbox{resume}{colback=orange!6,colframe=orange,title=Résumé en 3 phrases,fonttitle=\bfseries}

\newcommand{\ar}[1]{\textarabic{#1}}
\newcommand{\vrai}{\textcolor{vert}{\textbf{VRAI}}}
\newcommand{\faux}{\textcolor{rouge}{\textbf{FAUX}}}

\begin{document}

\section*{Partie 4a : Combien on produit ? (Quantité / Variété)}

\begin{carte}
Partie 3 \checkmark\ $\rightarrow$ \textbf{Partie 4a} $\rightarrow$ 4b (implantations) $\rightarrow$ 4c (MTS / MTO)

\begin{Arabic}
\textbf{الدور ديال هاد الجزء:} الدراجة مصممة. دابا خاصنا نعرفو أشمن نوع ديال المعمل غادي نحتاجو.
وأول سؤال هو: \textbf{شحال من دراجة غادي نصنعو؟ ونوع واحد ولا بزاف ديال الأنواع؟}
\end{Arabic}
\end{carte}

\subsection*{1. La quantité (Volume)}
\begin{center}
\begin{tabular}{lll}
\toprule
\textbf{Type} & \textbf{Quantité} & \textbf{Exemple} \\
\midrule
Unitaire & 1 pièce & bateau \\
Petite série & $< 100$ & vélos de course spéciaux \\
Moyenne série & $100$ à $10\,000$ & \\
Grande série & $> 10\,000$ & vélos normaux \\
Production de masse & très grands volumes & bouteilles d'eau \\
\bottomrule
\end{tabular}
\end{center}

\subsection*{2. La variété}
\begin{Arabic}
التنوع هو \textbf{شحال من نوع} ديال المنتوج كتصنع (لون، حجم، موديل…).

\textbf{القاعدة الذهبية:}
\begin{itemize}
\item \textbf{كمية كبيرة} ← \textbf{تنوع قليل} (بزاف ديال الدراجات بنفس الموديل)
\item \textbf{كمية صغيرة} ← \textbf{تنوع كبير} (شوية ديال الدراجات، وكل وحدة مختلفة)
\end{itemize}
\end{Arabic}

\begin{examen}
« La variété est en \textbf{corrélation inverse} avec la quantité. »

« Moins la variété des produits est grande, plus la quantité est grande. » \vrai
\end{examen}

\subsection*{3. La répétitivité}
\begin{itemize}
\item \textbf{Répétitif} : même produit, en grande quantité.
\item \textbf{Non répétitif} : produits uniques ou petites séries, souvent personnalisés.
\end{itemize}

\subsection*{4. Productivité et flexibilité \quad (tombe à chaque examen)}
\begin{Arabic}
جوج كلمات خاصك تفهمهم مزيان:
\begin{itemize}
\item \textbf{\textfrench{Productivité}} (الإنتاجية) = تصنع \textbf{بزاف} بزربة ← مرتبطة بـ\textbf{الكمية}
\item \textbf{\textfrench{Flexibilité}} (المرونة) = تقدر \textbf{تبدل} المنتوج بسهولة ← مرتبطة بـ\textbf{التنوع}
\end{itemize}
ما يمكنش يكون عندك جوج بزاف فنفس الوقت.
\end{Arabic}

\begin{examen}
« Plus on gagne en productivité, plus on perd en flexibilité. »
\end{examen}

\begin{center}
\begin{tabularx}{\linewidth}{l >{\raggedleft\arraybackslash}X c c}
\toprule
\textbf{Système} & \textbf{Description} & \textbf{Flexibilité} & \textbf{Productivité} \\
\midrule
Atelier classique & \ar{ماكينات عادية، كل مرة حاجة مختلفة} & élevée & faible \\
Atelier flexible (FMS) & \ar{ماكينات أوتوماتيكية كتقدر تبدل المنتوج} & très élevée & moyenne \\
Cellule flexible & \ar{مجموعة ماكينات لعائلة وحدة ديال القطع} & moyenne & moyenne / élevée \\
Machine spéciale & \ar{ماكينة لقطعة وحدة بكمية كبيرة} & faible & élevée \\
Ligne dédiée & \ar{خط إنتاج لمنتوج واحد بكمية ضخمة} & très faible & très élevée \\
\bottomrule
\end{tabularx}
\end{center}

\begin{piege}
\begin{itemize}
\item « Plus la variété est grande, plus la quantité est grande. » \faux
\item « Le nombre de différents types de produits quantifie la productivité. » \faux\ \ar{(التنوع = \textfrench{flexibilité})}
\item « La quantité de produits quantifie la flexibilité. » \faux\ \ar{(الكمية = \textfrench{productivité})}
\item « On choisit le type de système de production selon les deux critères productivité et flexibilité. » \vrai
\end{itemize}
\end{piege}

\begin{resume}
\begin{enumerate}
\item Quantité grande $\Rightarrow$ variété faible (corrélation inverse).
\item Quantité $\rightarrow$ productivité \quad / \quad Variété $\rightarrow$ flexibilité.
\item Plus on gagne en productivité, plus on perd en flexibilité.
\end{enumerate}
\end{resume}

\subsection*{Mini-QCM (réponses en bas)}
\begin{enumerate}
\item Relation entre variété et quantité ? \quad (a) plus de variété = plus de quantité \quad (b) moins de variété = plus de quantité \quad (c) indépendantes
\item Le nombre de types de produits quantifie : \quad (a) la productivité \quad (b) la flexibilité \quad (c) la qualité
\item Le système le plus productif : \quad (a) atelier classique \quad (b) cellule flexible \quad (c) ligne dédiée
\item Une moyenne série : \quad (a) $<100$ \quad (b) $100$ à $10\,000$ \quad (c) $>10\,000$
\end{enumerate}
{\small\textit{Réponses : 1-b, 2-b, 3-c, 4-b.}}

\end{document}
